\documentclass[11pt,a4paper]{article}
\usepackage{curso}
\begin{document}
\portada{Ejercicios --- Lección}{7.6}{Realización física con redes RC}
\section{Ejercicios}
\begin{enumerate}[label=\textbf{E\arabic*.}]
\item (B) Deduzca la lectura de equilibrio $C_e/\sum_xC_x$ con conservación de carga.
\item (B) ¿Cuántos bits para distinguir órdenes hasta $B=10^6$?
\item (B) Tolerancia de desajuste para $r=1000$.
\item (B) ¿Por qué hace falta una dirección por residuo?
\item (I) Con la cota $e^{-2gt/C_0}$, ¿qué tiempo garantiza error $\le1/(8B^2)$? Evalúelo para $B=10^6$.
\item (I) Si el desajuste es uniforme ($\varepsilon_x=\varepsilon$ para todo $x$), ¿qué ocurre con la lectura?
\item (A) Simule la red de $N=299$ residuos con desajuste $\sigma$ y ruido $k_BT/C$, y encuentre la región $(\sigma,J)$ donde la tasa de acierto supera el 99\,\%.
\end{enumerate}
\section{Soluciones}
\begin{solucion}
\textbf{E1.} $\sum_xC_xV_x$ se conserva y los voltajes tienden a un valor común $V_\infty$: $V_\infty\sum C=C_eV_0$.

\textbf{E2.} $2\log_2(10^6)\approx39{,}9$, unos 40 bits más bits de guarda.

\textbf{E3.} $|\bar\varepsilon-\varepsilon_e|<1/2000=0{,}05\,\%$.

\textbf{E4.} Un ciclo de longitud $r$ presupone $r$; con $N$ direcciones la carga se queda en la órbita de tamaño desconocido.

\textbf{E5.} $t\ge\frac{C_0}{2g}\ln(8B^2)$; para $B=10^6$: $\approx\frac{C_0}{2g}\cdot29{,}7\approx14{,}9\,C_0/g$.

\textbf{E6.} Se cancela: $(1+\varepsilon)/(1+\varepsilon)=1$; la lectura es exactamente $r$.

\textbf{E7.} Hace falta $\sigma\lesssim1/(2r)$ \emph{y} $J$ suficiente para que el ruido rms quede por debajo de $1/(8r^2)$.
\end{solucion}
\end{document}
